\documentclass[11pt]{article}
\usepackage{icmlko}
\begin{document}

\title{어려운 행은 무리가 아니었다 --- 일곱 도메인에서 난이도 변수를
못 찾았다}
\author{월드모델 실험실}
\date{노트 483 · 2026년 8월 3일}
\maketitle

\begin{abstract}
``어려운 행''이 두 노트째 나왔다 --- 노트 480 은 \emph{축 많은 행이 쉽다},
노트 482 는 \emph{\texttt{venue\_prominence} 붙은 $119$행이 어렵다}.
둘 다 \textbf{설명을 행으로 미뤘다}. 그러면 물어야 한다 --- \textbf{그
어려움이 무엇인가.}

유보 행마다 순위 오차를 내고 다섯 변수와 견줬다. 일곱 도메인 전부.

\begin{center}
\begin{tabular}{lrl}
\hline
변수 & 평균 $|$rho$|$ & 부호 \\
\hline
라벨 값 & $0.096$ & \textbf{갈린다}(애니 $+0.200$ · 세계애니 $-0.219$) \\
\textbf{라벨 극단성} & $0.088$ & \textbf{$7/7$ 음수} \\
축 수 & $0.068$ & 대체로 음수 \\
연도 & $0.055$ & 갈린다 \\
\texttt{venue} 유무 & $0.042$ & 제일 약함 \\
\hline
\end{tabular}
\end{center}

등록한 문장(\emph{극단성이 제일 세다})은 크기로는 \textbf{반증}이다 ---
라벨 값이 $0.096$ 으로 더 크다. 그런데 \textbf{라벨 값은 부호가 도메인
마다 갈린다.} 애니에서는 큰 라벨이 어렵고 세계애니에서는 쉽다. 법칙이
아니다.

\textbf{부호가 일관된 것은 극단성 하나뿐이고, 전부 $|$rho$| \le 0.22$ 로
작다.}
\end{abstract}

\begin{figure}[t]
\centering
\includegraphics[width=\textwidth]{figs/f1.pdf}
\caption{왼쪽 --- 일곱 도메인 $\times$ 다섯 변수. 초록 테두리(극단성)만
한 방향이고 붉은 점선(라벨 값)은 위아래로 갈린다. 오른쪽 --- 크기는
라벨 값이 크지만 부호는 극단성만 모인다.}
\end{figure}

\section{극단성은 지표의 성질이지 발견이 아니다}

$7/7$ 음수 --- \textbf{라벨이 가운데 몰린 행일수록 오차가 크다}. 이건
어느 정도 당연하다. 순위 지표에서 가운데 있는 행들은 서로 라벨 차이가
작아 원래 안 갈린다. \emph{모형이 못해서가 아니라 잴 것이 적어서}다.

그리고 크기가 $0.088$ 이다. \textbf{그것으로 노트 482 의 $-0.098$ 짜리
격차를 못 만든다.}

\section{그래서 ``어려운 행''은 우리가 가진 것으로는 못 짚는다}

이 노트의 결론은 부정이다.

\begin{quote}
\textbf{노트 $480 \cdot 482$ 가 ``행이 원래 어렵다''로 미룬 것을, 우리가
잰 어떤 변수로도 짚을 수 없다.} 라벨 값 · 극단성 · 연도 · 축 수 ·
\texttt{venue} 유무 다섯 중 어느 것도 부호가 일관되면서 크기가 있지 않다.
\end{quote}

\textbf{그 어려움은 미측정 요인이다.} 자료에 없는 무언가가 그 $119$행을
어렵게 만든다 --- 그것이 무엇인지는 이 자료로 못 말한다.

\section*{조항 --- ``행이 어렵다''는 설명이 아니다}

\begin{quote}
\textbf{조항} --- 어떤 격차를 \emph{행 난이도}로 설명했으면, \textbf{그
난이도를 재는 변수를 대야 한다.} 못 대면 그 설명은 \emph{이름 붙이기}이지
설명이 아니다. 노트 $480 \cdot 482$ 의 ``원래 어려운 행''은 지금 그
상태다 --- \textbf{맞는 이름이지만 아직 설명은 아니다.}
\end{quote}

이것은 노트 133(\emph{첫 양수를 그대로 채택하지 않는다})의 사촌이다.
그때는 \emph{수}를 의심했고 여기서는 \emph{설명}을 의심한다.

\section*{남은 것}

\begin{itemize}
\item \textbf{미측정 요인을 찾으려면 자료를 열어야 한다} --- 그
\texttt{venue\_prominence} 가 붙는 애니 $119$편이 어떤 작품들인지.
모형 갈래가 아니라 수집 갈래다.
\item \textbf{극단성 $0.088$ 은 채점 방식의 성질} --- 가운데 행을 덜
세는 가중을 쓰면 판이 달라질까. 순위 규약을 안 건드리기로 했으니
기록만 한다.
\item \textbf{라벨 값의 부호가 왜 도메인마다 갈리나} --- 애니 $+0.200$
대 세계애니 $-0.219$ 는 같은 애니메이션인데 반대다.
\end{itemize}

\end{document}
